\begin{table}[t]
\centering
\caption{Exactitud estandarizada (SA promedio de 100 sorteos) con $n$ proyectos de entrenamiento, en el formato log-lineal/bosque aleatorio. La prueba usa todos los dem\'as proyectos; la celda queda vac\'{\i}a cuando $n$ no es menor que la mitad de los $N$ proyectos del dataset.}
\label{tab:curva}
\scriptsize
\setlength{\tabcolsep}{2.5pt}
\resizebox{\ifdim\width>\columnwidth\columnwidth\else\width\fi}{!}{%
\begin{tabular}{lrcccc}
\toprule
Dataset & $N$ & $n=4$ & $n=8$ & $n=16$ & $n=32$ \\
\midrule
Albrecht & 24 & 0.29/0.31 & 0.66/0.49 &  &  \\
Maxwell & 62 & \ensuremath{-}0.51/0.23 & 0.46/0.35 & 0.53/0.42 &  \\
COCOMO81 & 63 & \ensuremath{-}2.67/0.29 & 0.30/0.33 & 0.50/0.46 &  \\
Desharnais & 81 & \ensuremath{-}0.34/0.27 & 0.40/0.38 & 0.49/0.42 & 0.51/0.45 \\
NASA93 & 93 & \ensuremath{-}0.10/0.25 & 0.43/0.38 & 0.50/0.46 & 0.56/0.56 \\
Kitchenham & 145 & \ensuremath{-}6.02/0.26 & 0.42/0.33 & 0.46/0.37 & 0.50/0.43 \\
China & 499 & \ensuremath{-}62.35/0.26 & 0.28/0.34 & 0.45/0.40 & 0.50/0.44 \\
\midrule
Mediana & & \ensuremath{-}0.51/0.26 & 0.42/0.35 & 0.49/0.42 & 0.50/0.45 \\
\bottomrule
\end{tabular}%
}
\par\smallskip
{\scriptsize\raggedright Con $n=4$ el promedio del log-lineal lo arrastran pocos sorteos en que la recta de cuatro puntos da pendientes extremas; en el sorteo t\'{\i}pico (la mediana de los 100) el log-lineal iguala o supera al bosque en 7 de 7 datasets.\par}
\end{table}
